
%-------------------------
% Resume in LaTeX
% Author: Azizzhan Supiyev (adapted from Jake's Resume template)
%-------------------------
\documentclass[letterpaper,11pt]{article}

\usepackage{latexsym}
\usepackage[empty]{fullpage}
\usepackage{titlesec}
\usepackage{marvosym}
\usepackage[usenames,dvipsnames]{color}
\usepackage[dvipsnames]{xcolor}
\usepackage{verbatim}
\usepackage{enumitem}
\usepackage[hidelinks]{hyperref}
\usepackage{fancyhdr}
\usepackage[italian,english]{babel}
\usepackage{tabularx}
\input{glyphtounicode}

\pagestyle{fancy}
\fancyhf{}
\fancyfoot{}
\renewcommand{\headrulewidth}{0pt}
\renewcommand{\footrulewidth}{0pt}

\addtolength{\oddsidemargin}{-0.5in}
\addtolength{\evensidemargin}{-0.5in}
\addtolength{\textwidth}{1in}
\addtolength{\topmargin}{-.5in}
\addtolength{\textheight}{1.0in}

\urlstyle{same}
\hypersetup{
  colorlinks=true,      % paint the text (instead of putting boxes)
  urlcolor=NavyBlue,    % colour for \href and URLs
  linkcolor=NavyBlue,   % internal links (section refs, etc.)
  citecolor=NavyBlue    % bibliography links
}

\raggedbottom
\raggedright
\setlength{\tabcolsep}{0in}
\setlength{\footskip}{4.08003pt}

\titleformat{\section}{
  \vspace{-4pt}\scshape\raggedright\large
}{}{0em}{}[\color{black}\titlerule \vspace{-5pt}]

\pdfgentounicode=1

\newcommand{\resumeItem}[1]{\item\small{{#1 \vspace{-2pt}}}}
\newcommand{\resumeSubheading}[4]{
  \vspace{-2pt}\item
    \begin{tabular*}{0.97\textwidth}[t]{l@{\extracolsep{\fill}}r}
      \textbf{#1} & #2 \\
      \textit{\small#3} & \textit{\small #4} \\
    \end{tabular*}\vspace{-7pt}}
\newcommand{\resumeSubSubheading}[2]{
    \item
    \begin{tabular*}{0.97\textwidth}{l@{\extracolsep{\fill}}r}
      \textit{\small#1} & \textit{\small #2} \\
    \end{tabular*}\vspace{-7pt}}
\newcommand{\resumeProjectHeading}[2]{
    \item
    \begin{tabular*}{0.97\textwidth}{l@{\extracolsep{\fill}}r}
      \small#1 & #2 \\
    \end{tabular*}\vspace{-7pt}}
\newcommand{\resumeSubItem}[1]{\resumeItem{#1}\vspace{-4pt}}
\renewcommand\labelitemii{$\vcenter{\hbox{\tiny$\bullet$}}$}
\newcommand{\resumeSubHeadingListStart}{\begin{itemize}[leftmargin=0.15in, label={}]}
\newcommand{\resumeSubHeadingListEnd}{\end{itemize}}
\newcommand{\resumeItemListStart}{\begin{itemize}}
\newcommand{\resumeItemListEnd}{\end{itemize}\vspace{-5pt}}

\begin{document}
\selectlanguage{italian}

%----------HEADING----------
\begin{center}
    \textbf{\Huge \scshape Azizzhan Supiyev} \\ \vspace{1pt}
    \small Software Engineer \\
    \href{mailto:azeasy19@gmail.com}{azeasy19@gmail.com} $|$
    \href{https://linkedin.com/in/azeasy}{LinkedIn} $|$
    \href{https://github.com/azeasy}{GitHub} $|$
    Telegram: \href{https://t.me/azeasy}{@azeasy}
\end{center}

%-----------SUMMARY-----------
\section{Sommario}
Ingegnere software esperto specializzato in \textbf{Python/Django}, API REST e architetture scalabili. Competenze comprovate nell'integrazione di sistemi complessi (OpenAI, Google, LinkedIn, API Gekata), nella progettazione di pipeline di dati e nell'applicazione di un solido ragionamento matematico per fornire soluzioni robuste.

%-----------EXPERIENCE-----------
\section{Esperienza}
\resumeSubHeadingListStart

\resumeSubheading
{Akhter Studios}{novembre 2023 - febbraio 2026}
{Ingegnere informatico}{Almaty, KZ}
\resumeItemListStart
    \resumeItem{\textbf{Introdex}: integrazione di Google, LinkedIn e altre API; ricerca ottimizzata tramite normalizzazione e indicizzazione del database legacy, velocità di ricerca ridotta da 5-7 secondi a meno di 300 ms.}
    \resumeItem{\textbf{MyStudy}: Sviluppo di una piattaforma educativa e di un motore di generazione di test basato sull'intelligenza artificiale, ottimizzazione dei prompt per il modello a basso costo al fine di ridurre le allucinazioni e aumentare la precisione.}
    \resumeItem{\textbf{AI Vendor Discovery Agent}: Progettazione di un chatbot basato su Django utilizzando Llama 3 con funzionalità personalizzate di richiamo degli strumenti per automatizzare l'abbinamento dei fornitori e le query complesse del database tramite Telegram o WhatsApp.}
    \resumeItem{Tecnologia: Django, DRF, Asyncio, Celery, PostgreSQL, Redis, Docker, Pytest, FastAPI, LangChain, LangGraph, LLM (OpenAI, xAI, LLama)}
\resumeItemListEnd

\resumeSubheading
{Adal-Meteo}{Gennaio 2020 - Ottobre 2023}
{Ingegnere informatico}{Almaty, KZ}
\resumeItemListStart
    \resumeItem{Sviluppo di API REST, pipeline ETL e dashboard in tempo reale per il monitoraggio ecologico.}
    \resumeItem{Implementato un servizio ETL con dashboard visivi per il confronto dei sistemi di previsione.}
    \resumeItem{Ha gestito i server locali e fornito assistenza al portale meteorologico governativo.}
    \resumeItem{Tecnologia: Django, Flask, server TCP, PostgreSQL, Docker, Python Dash.}
\resumeItemListEnd

\resumeSubheading
{Alem Research}{Gennaio 2019 - Gennaio 2020}
{Ingegnere informatico}{Almaty, KZ}
\resumeItemListStart
    \resumeItem{Progettazione di una pipeline ETL microservizio per il monitoraggio dei media utilizzando RabbitMQ, Redis, Solr.}
\resumeItemListEnd

\resumeSubHeadingListEnd

%-----------EDUCATION-----------
\section{Istruzione}
\resumeSubHeadingListStart
\resumeSubheading
{Università di Padova}{Ottobre 2022 - Luglio 2023}
{Laurea Magistrale in Matematica (Algebra e Geometria)}{Padova, Italia}
\resumeSubheading
{Università tecnica kazako-britannica}{agosto 2017 - giugno 2021}
{Laurea in Matematica Applicata}{Almaty, Kazakistan}
\resumeSubHeadingListEnd

%-----------SKILLS-----------
\section{Competenze tecniche}
\resumeSubHeadingListStart
\resumeItem{Lingue: Python (esperto), Golang (conoscenza di base)}
\resumeItem{Framework: Django, DRF, Flask, }
\resumeItem{Database: PostgreSQL, Redis, MongoDB}
\resumeItem{Strumenti: Celery, Asyncio, Docker, Git, RabbitMQ, Pytest}
\resumeItem{Altro: API REST, pipeline ETL, server TCP, NumPy, pandas}
\resumeSubHeadingListEnd

%-----------CERTIFICATIONS & TRAINING.-----------
\section{Certificazioni e formazione}
\resumeSubHeadingListStart
\resumeItem{Programma di formazione Epam Data Engineering: Databricks, PySpark, pipeline ETL}
\resumeItem{Percorso Ozon256: Golang, grpc, microservizi, Kafka, concorrenza, osservabilità}
\resumeSubHeadingListEnd

%-----------LANGUAGES-----------
\section{Lingue}
\resumeSubHeadingListStart
\resumeItem{Inglese (C1), Russo (Fluente), Kazako (B2), Italiano (B1)}
\resumeSubHeadingListEnd

\end{document}
